
\medskip
\phantomsection\label{lem:sga3-VIB-5.6.1}%
\noindent\textbf{Lemma 5.6.1.---}
\emph{Let \(S\) be a scheme, let \(f:X\to S\) be an \(S\)-scheme
locally of finite presentation, let \(s\) be a point of \(S\), and
let \(x\) be a closed point of \(X_s\). Suppose that \(f\) is
universally open at the generic point \(z\) of an irreducible
component \(Z\) of \(X_s\) containing \(x\), and that}
\[
  \dim_x(X_s)=\dim Z.
\]
\emph{Then there exists a commutative diagram}
\begin{center}
  \phantomsection\label{diag:sga3-VIB-5.6.1-statement}%
  \begin{tikzcd}[row sep=large,column sep=large]
    X \arrow[r,"f"] &
    S \\
    S''
      \arrow[u,"h"]
      \arrow[r,"\pi"]
      \arrow[ur,"\phi"] &
    S' \arrow[u,"w"']
  \end{tikzcd}\enspace,
\end{center}
\emph{where \(S'\) is affine, \(w\) is \'etale, \(\pi\) is finite,
surjective, and of finite presentation, and
\(\phi^{-1}(s)\) consists of a single point \(s''\), such that
\(h(s'')=x\) and \(\kappa(s'')=\kappa(x)\).}%
\footnote{\label{note:sga3-VIB-41}Editor's note: consequently,
\(h\) induces a section \(\sigma\) of
\(X\times_S S''\to S''\) such that \(\sigma(s'')\) lies above
\(x\); compare EGA~IV\(_3\), 14.5.10.}

\emph{Proof.}
First, we may suppose that \(S=\operatorname{Spec}(A)\) and
\(X=\operatorname{Spec}(B)\), where \(B\) is an \(A\)-algebra of
finite presentation. Let \(\mathfrak p\) and \(\mathfrak q\) be
the prime ideals of \(A\) and \(B\) corresponding respectively to
\(s\) and \(x\), so that
\[
  \mathfrak p=A\cap\mathfrak q.
\]
Set \(n=\dim_x(X_s)\), and choose \(t_1,\ldots,t_n\in B\) whose
images in
\[
  \mathcal O_{X_s,x}=B_{\mathfrak q}/\mathfrak p B_{\mathfrak q}
\]
form a system of parameters. Then
\[
  \mathcal O_{X_s,x}/(t_1,\ldots,t_n)
\]
is finite-dimensional over \(\kappa(x)\), and hence also over
\(\kappa(s)\), because \(x\) is a closed point of the algebraic
\(\kappa(s)\)-scheme \(X_s\). Consequently, the \(S\)-morphism
\[
  u:X\longrightarrow\mathbb A^n_S
    =\operatorname{Spec}\bigl(A[T_1,\ldots,T_n]\bigr),
  \qquad T_i\longmapsto t_i,
\]
is of finite presentation, \(x\) is isolated in the fiber
\(u^{-1}(u(x))\), and
\[
  u(x)=\tau_0(s),
\]
where \(\tau_0:S\to\mathbb A^n_S\) denotes the zero section of
\(\mathbb A^n_S\to S\), corresponding to the \(A\)-algebra
homomorphism
\[
  A[T_1,\ldots,T_n]\longrightarrow A,\qquad T_i\longmapsto0.
\]
The set of points of \(X\) that are isolated in their fibers over
\(\mathbb A^n_S\) is open (EGA~IV\(_3\), 13.1.4). Thus, after
shrinking \(X\), we may suppose that \(u\) is quasi-finite and
that \(u^{-1}(u(x))=\{x\}\). By \hyperref[lem:sga3-VIB-5.6.1.0]{Lemma 5.6.1.0}, \(u\) is
universally open at \(x\).

Let
\[
  B_0=B/(t_1,\ldots,t_n),\qquad
  X_0=X\times_{\mathbb A^n_S}\tau_0(S)
      =\operatorname{Spec}(B_0),
\]
and let \(u_0:X_0\to S\) be obtained from \(u\) by the base change
\(\tau_0:S\hookrightarrow\mathbb A^n_S\). Then \(u_0\) is
quasi-finite and of finite presentation, is universally open at
\(x\), and \(x\) is the unique point of \(X_0\) above \(s\).

Let \(A'\) be the henselization of the local ring
\[
  A_{\mathfrak p}=\mathcal O_{S,s},
\]
and set
\[
  S'=\operatorname{Spec}(A'),\qquad
  B'_0=B_0\otimes_A A',\qquad
  X'_0=X_0\times_S S'=\operatorname{Spec}(B'_0).
\]
The closed point \(s'\) of \(S'\) is the unique point of \(S'\)
above \(s\), and \(\kappa(s')=\kappa(s)\). The scheme \(X'_0\)
has a unique point \(x'\) above \(s'\); moreover
\(\kappa(x')=\kappa(x)\), and \(x'\) is also the unique point of
\(X'_0\) above both \(s\) and \(x\). Since \(A'\) is henselian,
EGA~IV\(_4\), 18.5.11 gives a disjoint decomposition into open and
closed parts
\[
  \tag{*}\label{sga3:vi:lemma:5-6-1:equation:star}
  X'_0=V\amalg W,
  \qquad
  V=\operatorname{Spec}(\mathcal O_{X'_0,x'}).
\]
The local ring \(\mathcal O_{X'_0,x'}\) is finite and of finite
presentation over \(A'\). Hence the restriction \(\pi\) of
\(u'_0\) to \(V\) is finite and of finite presentation. Moreover,
because \(u'_0:X'_0\to S'\) is open at \(x'\),
\[
  u'_0(V)=\operatorname{Spec}(\mathcal O_{S',s'})=S',
\]
so \(\pi\) is surjective.

This proves the assertion when \(S=S'\). In general, \(A'\) is
the filtered direct limit of \(A\)-subalgebras \(A_i\) that are
\'etale over \(A\) and for which
\[
  S_i=\operatorname{Spec}(A_i)
\]
has a unique point \(s_i\) above \(s\), with
\(\kappa(s_i)=\kappa(s)\). We have
\[
  B'_0=\varinjlim_i B_i,
  \qquad
  B_i=B_0\otimes_A A_i.
\]
Set
\[
  X_i=\operatorname{Spec}(B_i)=X_0\times_S S_i,
  \qquad
  C=\mathcal O_{X'_0,x'}.
\]
By \((*)\), one has
\[
  C\simeq B'_0/fB'_0
\]
for some idempotent \(f\in B'_0\). There are an index \(i\) and
an element \(f_i\in B_i\) whose image in \(B'_0\) is \(f\). Put
\[
  C_i=B_i/f_iB_i,\qquad V_i=\operatorname{Spec}(C_i).
\]
Then
\[
  C=C_i\otimes_{A_i}A',
  \qquad
  V=V_i\times_{S_i}S',
\]
and \(C_i\) is an \(A_i\)-algebra of finite presentation, since
the same is true of \(B_i\). Consequently, the morphisms
\begin{center}
  \phantomsection\label{diag:sga3-VIB-5.6.1-limit-source}%
  \begin{tikzcd}[row sep=large,column sep=large]
    X'_0 \arrow[dr,"u'_0"] &
    {} \\
    V
      \arrow[u,"h'"]
      \arrow[r,"\pi"] &
    S'
  \end{tikzcd}
\end{center}
arise, by the base change \(S'\to S_i\), from morphisms of finite
presentation
\begin{center}
  \phantomsection\label{diag:sga3-VIB-5.6.1-limit-descent}%
  \begin{tikzcd}[row sep=large,column sep=large]
    X_i \arrow[dr,"u_i"] &
    {} \\
    V_i
      \arrow[u,"h_i"]
      \arrow[r,"\pi_i"] &
    S_i
  \end{tikzcd}
\end{center}

For every \(j\geq i\), let
\[
  V_j=V_i\times_{S_i}S_j
\]
and let \(\pi_j:V_j\to S_j\) be the morphism of finite
presentation obtained from \(\pi_i\) by base change. Since
\(\pi:V\to S'\) is finite and surjective, EGA~IV\(_2\), 8.10.5
shows that, for some \(j\), the morphism
\(\pi_j:V_j\to S_j\) is finite and surjective. Then
\(w:S_j\to S\) is affine and \'etale, \(S_j\) has a unique point
\(s_j\) above \(s\), and \(V_j\) has a unique point \(x_j\) above
\(s_j\), since \(x'\) is the unique point of
\[
  V=V_j\times_{S_j}S'
\]
above \(s'\):
\begin{center}
  \phantomsection\label{diag:sga3-VIB-5.6.1-final-composite}%
  {\small
  \begin{tikzcd}[row sep=large,column sep=large]
    X_j
      \arrow[r]
      \arrow[dr,"u_j"] &
    X_0
      \arrow[r]
      \arrow[dr,"u_0"] &
    X
      \arrow[d,"f"']
      \arrow[dr,"u"] &
    {} \\
    V_j
      \arrow[u,"h_j"]
      \arrow[r,"\pi_j"] &
    S_j
      \arrow[r,"w"] &
    S
      \arrow[r,"\tau_0"] &
    \mathbb A^n_S
  \end{tikzcd}}
\end{center}
Thus \(x_j\) is the unique point of \(V_j\) above \(s\), its image
under \(V_j\to X_j\to X\) is \(x\), and
\[
  \kappa(x_j)=\kappa(x).
\]
